%% kapa: exam=2 round=37 year=2026 period=2 subject=appraisal_theory subject_ko=감정평가이론 seq=2 number=2 points=30
%% source: Q-Net 감정평가사 기출문제 (cst00309) artlSeq=5268128 file=raw/1008/5268128/26년 제37회 감정평가사 2차 (2교시_감정평가이론).pdf
%% conversion: kordoc 4.24.51 -> tools/gen2.py (auto); review=proofread-gpt-terra; flags=none
\begin{longquestion}{【문제 2】}
감정평가활동에서 감정평가액의 변동성이 실제 거래가격의 변동성보다 작게 나타나거나, 시장의 실제 변화속도보다 느리게 나타나는 현상이 있다. 다음 물음에 답하시오. (30점) \par
\medskip\noindent 물음 1) 이러한 현상의 원인을 부동산시장의 특성 및 감정평가사의 의사결정과 관련하여 설명하시오. (단, 관계 법령 및 윤리규정 등의 위반은 없다고 가정함) (15점) \par
\medskip\noindent 물음 2) A시의 특정 아파트단지에서 최근 거래가격이 급격히 하락했으나 기준시점이 최근인 감정평가액은 이전 수준을 유지하고 있다. 이러한 감정평가액이 감정평가의 신뢰성에 미치는 영향을 시장가치와 관련하여 설명하시오. (15점) \par
\end{longquestion}
